%Preamble
\documentclass[12pt,a4paper]{article}

\usepackage[T1]{fontenc}
\usepackage[utf8]{inputenc}
\usepackage[romanian]{babel}
\usepackage{lmodern}
\usepackage[margin=2.5cm]{geometry}
\usepackage{microtype}
\usepackage{hyperref}

\setlength{\parindent}{0pt}
\setlength{\parskip}{0.8em}

\begin{document}

\textbf{Titlu:} Agent inteligent pentru medii virtuale interactive

\textbf{Student:} Ing.\ Marius-Vlad DUMITRU

\textbf{Descriere a temei:} Proiectul explorează dezvoltarea și integrarea agenților inteligenți autonomi în medii virtuale, cu accent pe percepție, luarea deciziilor și acțiune în conformitate cu un set de reguli predefinite. Acest demers contribuie la avansarea inteligenței artificiale, în special în domeniul luării deciziilor în timp real și al învățării adaptive. Prin analiza comportamentului unui agent inteligent într-un mediu virtual, proiectul își propune să abordeze limitările actuale ale modelelor de învățare profundă, cum ar fi interpretabilitatea și explicabilitatea deciziilor luate. Rezultatele vor fi aplicate în cadrul unor jocuri video, cu scopul de a demonstra aplicabilitatea și eficiența agenților în scenarii interactive. În procesul de dezvoltare, se vor utiliza tehnologii și algoritmi moderni, implementați în Python și integrați cu platforme de simulare virtuală. Proiectul va include, de asemenea, o evaluare comparativă a performanței agenților în raport cu alte sisteme similare.

\textbf{Contribuția studentului:} Studentul va implementa modulele necesare pentru integrarea agenților inteligenți autonomi într-un mediu virtual interactiv, asigurând compatibilitatea cu platformele de simulare utilizate. Vor fi dezvoltate interfețe de interacțiune între agent și mediu, iar arhitectura acestora va fi ajustată pentru a îmbunătăți performanța sau pentru a răspunde cerințelor specifice ale scenariilor analizate. Studentul va antrena modele pentru agenții autonomi folosind seturi de date relevante, iar în cazul în care acestea nu sunt suficiente, va explora metode pentru generarea sau extinderea lor. De asemenea, studentul va efectua experimente prin modificarea unor componente de arhitectură sau practici de antrenare pentru a evalua impactul acestora asupra performanței agenților, incluzând timpii de decizie și rata de succes în sarcini complexe. Validarea rezultatelor obținute va include o comparație cu alte metode existente din literatura de specialitate, iar eficiența sistemului propus va fi evaluată prin metrici specifice pentru performanța agenților în medii virtuale. Această analiză va fi completată cu propuneri pentru îmbunătățirea tehnicilor existente, în vederea creșterii interpretabilității și explicabilității deciziilor luate de agent.

\textbf{Resurse și bibliografie:}

\begin{thebibliography}{9}
\setlength{\itemsep}{0.5em}

\bibitem{liu2024} R.-Z. Liu, Y. Shen, Y. Yu and T. Lu, ``Revisiting of AlphaStar,'' in \textit{IEEE Transactions on Games}, vol.~16, no.~2, pp.~317--330, June 2024, doi: 10.1109/TG.2023.3265975.

\bibitem{tao2023} J. Tao et al., ``Explainable AI for Cheating Detection and Churn Prediction in Online Games,'' in \textit{IEEE Transactions on Games}, vol.~15, no.~2, pp.~242--251, June 2023, doi: 10.1109/TG.2022.3173399.

\bibitem{singh2024} J. Singh, S. Rani and G. Srilakshmi, ``Towards Explainable AI: Interpretable Models for Complex Decision-making,'' 2024 International Conference on Knowledge Engineering and Communication Systems (ICKECS), Chikkaballapur, India, 2024, pp.~1--5, doi: 10.1109/ICKECS61492.2024.10616500.

\bibitem{li2024} W. Li, W. Liu, S. Shao, S. Huang and A. Song, ``Attention-Based Intrinsic Reward Mixing Network for Credit Assignment in Multiagent Reinforcement Learning,'' in \textit{IEEE Transactions on Games}, vol.~16, no.~2, pp.~270--281, June 2024, doi: 10.1109/TG.2023.3263013.

\end{thebibliography}

\end{document}
